\begin{frame}{Jagtretten – grundregler (§13–§15)}
    \begin{itemize}
        \item \textbf{§13}: Ingen må jage uden jagtret – medmindre samtykke fra jagtberettiget
        \item \textbf{§14}: Jagtretten tilkommer ejendommens ejer
          \begin{itemize}
            \item På ferske vande (uden særlig ejendomsret): Ejerne af tilstødende ejendomme (nærmeste del)
            \item Gælder ikke hvis ældre rettigheder giver retten til andre
          \end{itemize}
        \item \textbf{§15}: Fast bopæl i DK → ret til ikke-erhvervsmæssig jagt på fiskeriterritoriet (fra højeste vandstandslinje)
          \begin{itemize}
            \item Gælder ikke hvis ældre ret til tilgrænsende landarealer
            \item Undtagelse: EU/EØS-borgere med tilknytning til DK
          \end{itemize}
    \end{itemize}
\end{frame}